\section{Вступ}

\subsection{Мета та область застосування документа}

Цей документ є специфікацією програмних вимог (Software Requirements Specification, SRS) до інформаційної системи персонального обліку медіа-контенту. Мета документа~--- формально описати функціональні та нефункціональні вимоги до системи, визначити межі версії~v1 (MVP) та надати базу для розроблення проектної документації (SDD) і тестування. \sked{K10}

\subsection{Призначення системи}

Система є веб-застосунком, що дозволяє зареєстрованому користувачеві вести особистий каталог медіа-контенту~--- книг, серіалів і фільмів. Застосунок забезпечує облік переглянутого та прочитаного контенту: фіксацію статусів, оцінок, нотаток, тегів, повторних переглядів (Sessions), організацію елементів у добірки (Collections) і перегляд особистої статистики. \sked{K2,K4,K3,K5}

\subsection{Контекст і аудиторія документа}

Документ підготовлений у рамках навчального курсу «Проектування інформаційних систем» (SDD-цикл). Цільова аудиторія: викладач курсу та розробники системи. Читач цього документа повинен отримати достатньо інформації для розроблення архітектури та плану тестування без додаткових уточнень. \sked{K44}

\subsection{Структура документа}

\begin{itemize}
  \item \secref{sec:zahalnyi-opys}~--- загальний опис системи, межі MVP, ролі користувачів, операційне середовище.
  \item \secref{sec:hlosarii}~--- визначення ключових термінів.
  \item \secref{sec:funktsionalni-vymohy}~--- повний перелік функціональних вимог, згрупованих за функціональними блоками.
  \item \secref{sec:nefunktsionalni-vymohy}~--- нефункціональні вимоги: продуктивність, безпека, надійність, зручність, сумісність, валідація.
  \item \secref{sec:kryterii-pryiniattia}~--- критерії прийняття для кожної вимоги.
  \item \secref{sec:bdd}~--- BDD-сценарії у форматі Gherkin.
  \item \secref{sec:matrytsia}~--- матриця простежуваності вимог.
  \item Додаток~А~--- перелік функцій, що не входять до MVP.
\end{itemize}
